\chapter{Summary}
\label{sec:summary}

\section{Conclusion}

This thesis investigated multi-objective optimization methods for two
gripper-design problems: the placement of three suction cups and the design of
rigid parallel-jaw contours. Both were formulated as black-box problems and
implemented through a common interface that separates candidate generation from
problem-specific geometric evaluation.

NSGA-II, NSGA-III, MOEA/D, SMS-EMOA, and CMOPSO were evaluated on four
workpieces, with each gripper--workpiece combination treated as a separate
problem instance. Solution quality was assessed using IGD+ and cumulative
archive hypervolume (HV), and runtime was recorded as an indication of
computational cost. For the suction-cup problem, SMS-EMOA achieved the lowest
mean IGD+ on all four workpieces, and NSGA-II was also among the leading
methods. HV did not identify a clear leader on every workpiece. For the
parallel-jaw problem, NSGA-II achieved the highest mean HV on all four
workpieces, while the lowest mean IGD+ depended on the workpiece. MOEA/D
generally performed less well than the leading methods. SMS-EMOA and NSGA-II are
therefore promising choices for the investigated formulations and budgets, but no
universally best algorithm was established.

The Ax/BoTorch Bayesian runs (batch size 10, four workpieces in parallel) were
stopped after more than 90 hours, at approximately 450 trials per run, and
produced no final results. A subsequent minimal test on the parallel-jaw gripper
and the connecting rod (batch size 5, 496 evaluations, one run) reached an HV of
0.0281. This is about 87\,\% of the best evolutionary result (NSGA-II, 0.0323)
with about 5\,\% of the evaluations. However, model fitting and candidate
generation were very time-consuming, and repeating the optimization for several
runs, further workpieces, and the suction-cup gripper would multiply this cost.
Bayesian optimization therefore did not produce usable results in this work,
whereas the population-based methods completed all reported experiments.

\section{Limitations}
\label{sec:limitations}

The study used five independent runs per algorithm and workpiece for the
suction-cup problem and three for the parallel-jaw problem, so the statistics are
descriptive and small differences should be interpreted cautiously. Algorithm
parameters were not tuned individually, and some HV curves were still rising at
the end of the evaluation budget.

The true Pareto fronts are unknown. IGD+ was therefore computed against an
empirical reference front built from the final solutions of the compared runs,
and it measures performance relative to those runs. It also evaluates the final
solution sets, whereas the reported HV uses a cumulative archive, so the two
indicators do not assess exactly the same sets.

Runtimes are indicative only because the runs used GPUs with different computing
capacities. The seed values and hardware information were not stored in the
\texttt{results.json} metadata, so the GPU of each run cannot be identified and
the runs cannot be reproduced exactly from the stored results. The Bayesian
results are limited to a minimal single-run test ($n=1$) on one
gripper--workpiece pair, with a smaller batch size (5) than the stopped main
attempt (10) and about 5\,\% of the evaluation budget of the evolutionary
algorithms. They were compared only through HV, and their variability and their
behavior at larger budgets or on other workpieces are unknown. In the main
attempt, the four parallel processes shared computational resources, and the
logs did not separate model fitting, candidate generation, and evaluation times.
The causes of the long runtime therefore remain unresolved, and the outcome
cannot be generalized to Bayesian optimization as a whole.

The suction-cup parameterization does not explicitly enforce a minimum
distance between cup centers. Consequently, some candidate layouts may
contain overlapping cups and may not be physically realizable. The
radial bounds alone do not prevent overlap because the angular positions
are optimized independently.

The evaluation is entirely computational. The suction-cup seal model and
the voxel-based proximity and collision calculations of the parallel-jaw
evaluator do not establish physical grasp performance, and the optimized
configurations were not validated on real grippers. Manufacturing effects,
material behavior, and contact forces were not assessed. Only four
workpieces and restricted design spaces were investigated.

\section{Future Work}

The Bayesian comparison should be repeated under controlled hardware conditions
with separate timing of model fitting, candidate generation, and geometric
evaluation. The minimal test on the connecting rod can serve as a starting point
and should be extended to several independent runs, the remaining workpieces, and
the suction-cup gripper. Settings such as objective thresholds and batch size
should be investigated. Running the processes sequentially and saving
intermediate results would help locate bottlenecks and preserve partial runs.
Comparisons should consider solution quality at common evaluation counts as well
as wall-clock time.

Future experiments should store the metadata needed for reproducibility in the
result files, including the random seed of every run along with the GPU and CPU model. This would allow runs
to be repeated and runtimes to be compared across hardware.

Further algorithm studies could use more repetitions, statistical tests, and
systematic parameter studies. Hybrid and surrogate-assisted methods could reduce
the number of expensive evaluations. Since the Bayesian run reached a high HV
after few evaluations, a short Bayesian phase could, for example, initialize an
evolutionary algorithm.

A non-overlap condition should be added to the suction-cup formulation, requiring
the distance between every pair of cup centers to be at least the cup diameter,
and existing layouts should be checked for overlap. The evaluations could also be
extended with additional design variables and force-based criteria, and the
voxel-based parallel-jaw calculations should be compared with direct collision
tests. Finally, physical experiments with selected suction-cup layouts and jaw
contours in a bin-picking setup could show whether better computational
objectives lead to reliable grasping.